\begin{titlepage}
   \begin{center}          
   \includegraphics[scale=1.3]{Portada/Escudo_UNAM_eps}
     \large
    %   \vspace*{0.2cm}  
        \begin{center}   
      {\fontfamily{qag}\selectfont
  \textbf{	UNIVERSIDAD NACIONAL AUTÓNOMA DE MÉXICO}\\
     \textbf{ {\color{Blue} 	Posgrado en Ciencias Físicas}}\\
} 
\end{center}  
   \end{center}


        \vspace*{0.6cm}
        \begin{center}
        
      \Huge       
          {\fontfamily{qag}\selectfont
        \textbf{Aspectos de la Teoría Cuántica de los supercampos}         
              }          \\
        \vspace{0.8cm}  
        \large
            \textit{ Tesis que para optar por el grado de:}\\
            \vspace{0.15cm}
       {\color{blue}DOCTOR EN CIENCIAS (FÍSICA)}\\
        \vspace{0.8cm} 
      
       \textit{ Presenta:}      \\  
        \vspace{0.25cm}
        \Large
           {\fontfamily{qag}\selectfont
       {Enrique Jiménez Ramos}\\
        }
         \vspace{1.0cm}   
         \end{center}       
         \begin{large}         
         {\color{Blue}     
             
        \indent Dra. Myriam Mondragón Ceballos,  Instituto de Física, UNAM\\   
        % \vspace{0.2cm}    
        \indent Dr. {Jens Erler Weber},       Instituto de Física, UNAM\\  
         %  \vspace{0.2cm}       
         \indent Dr. Alfonso Mondragón Ballesteros,    Instituto de Física, UNAM\\ 
                  }                          
         \end{large}
             
        \begin{center}
        \normalsize
        CIUDAD DE MÉXICO, ABRIL 2016
        \end{center}
                 
 %   \end{center}
\end{titlepage}
